\chapter{Concluding Remarks}
\label{ch:concluding_remarks}

\section{Chapter Overview}
\label{sec:concl_overview}
In this chapter, a brief summary of the major achievements, technological breakthroughs, lessons learned and directions for future research concerning the proposed paper-based customizable virtual keyboard are provided. Using only one monocular web-camera and simple paper without employing any special technologies such as depth camera, Time-of-Flight camera, laser projector, wearable gadgets and even being performed in real time on an ordinary PC, the described system proves that ordinary paper can be turned into customizable virtual input devices of any resolution of an ordinary camera.

Review of the research process includes evaluation of achievement of the research objectives, technical issues and their solutions, gained knowledge and skills, practical applications and suggestions for future research. Section~\ref{sec:concl_objectives} deals with the fulfillment of each of the research objectives in comparison with literature review, software engineering and experimental testing results. Section~\ref{sec:concl_problems} provides information about the technical problems and their solutions. Section~\ref{sec:concl_reflection} discusses personal reflections about the technical background and the project experience. Section~\ref{sec:concl_business} focuses on practical implementations and business examples. And finally, Section~\ref{sec:concl_future} is dedicated to future research.

\section{Accomplishment of Research Objectives}
\label{sec:concl_objectives}
To validate the outcomes of this research, the approach of triangulation has been utilized. The triangulation approach connects the outcome of empirical testing done in Chapter~\ref{ch:testing_evaluation} with the aim of this research stated in Chapter~\ref{ch:introduction} via implementation modules and system requirements described in Chapters~\ref{ch:implementation} and~\ref{ch:system_requirements}.

\subsection{Triangulation Evaluation of Objective 1: Problem Identification and Literature Review}
\label{subsec:concl_obj1}
\begin{itemize}
    \item \textbf{Objective Statement:} \emph{To identify existing limitations in monocular vision-based virtual keyboards and synthesize a comprehensive literature taxonomy.}
    \item \textbf{Triangulation Evidence:} An ordered taxonomy and comparison table were generated in Chapter~\ref{ch:literature_review} through an analysis of 54 scientific articles. The research gaps found from the literature review include: (1) dependence on expensive devices such as depth sensor, ToF camera, active laser projector and powerful graphics processing unit, (2) ambiguity in monocular depth in the case where fingers are above flat surfaces, (3) dependence on the size of hands and distance between camera and hands, and (4) inability to adapt the layout dynamically on a regular piece of paper \cite{ReviewVirtualKeyboard2020, Ji2018Local, Posner2012Single}.
    \item \textbf{Validation Outcome:} \textbf{Objective Fully Accomplished.} The literature taxonomy has played a very significant role in defining the system architecture in Chapter~\ref{ch:system_requirements}.
\end{itemize}

\subsection{Triangulation Evaluation of Objective 2: Scale Normalization and Deep Learning Pipeline}
\label{subsec:concl_obj2}
\begin{itemize}
    \item \textbf{Objective Statement:} \emph{To design and implement a scale-invariant hand feature extraction pipeline and a temporal PyTorch deep learning classifier for touch contact detection.}
    \item \textbf{Triangulation Evidence:} Partition 2 Data Pipeline and Partition 3 Deep Learning Benchmarking Engine were used for Chapter~\ref{ch:implementation}. The normalization of scales of hand joints location was done based on the unit-less hand length $L_{\text{hand}} = \|\mathbf{P}_9 - \mathbf{P}_0\|_2$. Sliding windows consisted of 5 frames and overlapped by 2 frames ($\mathbf{X}_W \in \mathbb{R}^{5 \times 84}$). The sliding windows were used at 12~FPS rate. The scaled normalized features were used as an input to preserve deceleration information during contact. As can be seen from the benchmark results presented in Chapter~\ref{ch:testing_evaluation}, the F1-score of one-directional PyTorch LSTM model known as \texttt{best\allowbreak\_finger\allowbreak\_touch\allowbreak\_lstm.pth} is $0.963$ and has the fast CPU inference time of $2.08\text{ ms}$.
    \item \textbf{Validation Outcome:} \textbf{Objective Fully Accomplished.} The distance invariance test showed good results in the range of $30\text{ cm}$ to $100\text{ cm}$ ($F1 \ge 0.958$) and camera resolution invariance from 360p to 1080p ($F1 \in [0.961, 0.964]$) running in real time on standard commodity CPUs without a GPU.
\end{itemize}

\subsection{Triangulation Evaluation of Objective 3: System Suite Design and Engineering}
\label{sec:concl_obj3}
\begin{itemize}
    \item \textbf{Objective Statement:} \emph{To design and implement an integrated two-application software suite enabling arbitrary custom layout design and decoupled action combination multiplexing on the same printed paper layout.}
    \item \textbf{Triangulation Evidence:} Two applications were developed in Chapter~\ref{ch:implementation}: Application 1 (Layout Designer, \texttt{designer\_app.py}) and Application 2 (Runtime Engine). Application 1 is made up of PySide6 drag-and-drop interface, automatic border setting for the AprilTag, XML layout export, and vector PDF file generation with 300 DPI. Application 2 is made up of Setup GUI (Component A). This offers key ID configuration for keystrokes or shell commands, switching and saving action profiles that are kept in JSON files within the same printed sheet of paper without the need for printing again. Component B constitutes the background engine. It captures the video feed, AprilTag homography, PyTorch LSTM inference, and coordinate mapping ($P_{\text{XML}} = H \cdot P_{\text{pixel}}$).
    \item \textbf{Validation Outcome:} \textbf{Objective Fully Accomplished.} Chapter~\ref{ch:testing_evaluation} includes functional testing where all ten tests (FT-01 to FT-10) have been conducted successfully with a $100\%$ pass rate.
\end{itemize}

\subsection{Triangulation Evaluation of Objective 4: Empirical System Evaluation}
\label{subsec:concl_obj4}
\begin{itemize}
    \item \textbf{Objective Statement:} \emph{To conduct rigorous empirical benchmarking evaluating real-time pipeline latency, optical homography tracking, and human typing performance.}
    \item \textbf{Triangulation Evidence:} Benchmarking of the system was carried out in Chapter~\ref{ch:testing_evaluation} using CPU-based normal hardware. It was observed that processing time took $29.09\text{ ms}$ ($34.4\text{ FPS}$) when there was a minimal consumption of system resources in terms of $142.5\text{ MB}$ of RAM and $24.8\%$ of CPU utilization. This made the requirement of $83.33\text{ ms}$ of processing time for 12~FPS a very easy task. The AprilTag homography-based tracking RMS error was always less than $0.42\text{ mm}$ with a tilt of the camera at $60^\circ$. Users' typing test produced an average typing speed of $24.6\text{ WPM}$ and a Character Error Rate of $5.2\%$.
    \item \textbf{Validation Outcome:} \textbf{Objective Fully Accomplished.} This performance was better than previous virtual keyboards using a single camera ($12.8\text{ WPM}$, $11.4\text{ CER}$ \cite{Ji2018Local}).
\end{itemize}

Table~\ref{tab:objective_triangulation_summary} presents the objective accomplishment and triangulation matrix.

\begin{table}[htbp]
\centering
\caption{Research objective accomplishment and triangulation matrix.}
\label{tab:objective_triangulation_summary}
\resizebox{\textwidth}{!}{%
\begin{tabular}{p{3.5cm} p{3.5cm} p{4.0cm} p{3.5cm} c}
\toprule
\textbf{Research Objective} & \textbf{Design / Chapter Origin} & \textbf{Implementation Artifact} & \textbf{Empirical Validation Evidence} & \textbf{Status} \\
\midrule
\textbf{O1: Literature Review and Gap Identification} & Chapter 2 Literature Taxonomy & Research DB (\texttt{references.bib}, \texttt{summary-matrix.md}) & Identified 4 key gaps; established evaluation baseline. & \textbf{ACCOMPLISHED} \\
\midrule
\textbf{O2: Scale Normalization and Neural Touch Classifier} & Chapter 3 Methodology and Chapter 5 Pipeline & Partition 2 (\texttt{datacreator/}) and Partition 3 (\texttt{run\_all.py}) & F1-score $0.963$; CPU inference $2.08\text{ ms}$; CPU-only real-time at 12~FPS. & \textbf{ACCOMPLISHED} \\
\midrule
\textbf{O3: Two-Application System Suite Creation} & Chapter 4 Requirements and Chapter 5 Architecture & Application 1 (\texttt{designer\_app.py}) and Application 2 Suite & $100\%$ pass rate across FT-01--FT-10; same-layout multi-profile switching verified. & \textbf{ACCOMPLISHED} \\
\midrule
\textbf{O4: Empirical Latency and Typing Evaluation} & Chapter 6 Evaluation Framework & Multi-threaded Engine and Evaluation Scripts & Latency $29.09\text{ ms}$ ($34.4\text{ FPS}$); $24.6\text{ WPM}$; $5.2\text{ CER}$. & \textbf{ACCOMPLISHED} \\
\bottomrule
\end{tabular}%
}
\end{table}

\section{Problems Encountered and Technical Mitigations}
\label{sec:concl_problems}
A number of practical issues arose while developing the system due to constraints of optical cameras and human interaction.

\subsection{Hardware Constraints of Monocular RGB Capture}
As a rule, webcams provide videos at $30\text{ FPS}$ and have rolling shutter sensors that generate motion blur during finger taps with high speed \cite{Thomas2013Camera}. The high speed of typing means high speed of downward finger movements that can result in lower tracking accuracy of landmarks.
\begin{itemize}
    \item \textbf{Technical Mitigation:} During dataset creation, the PyTorch LSTM model was trained using resampling at 12~FPS (\texttt{resample\_12fps.py}). As a result, the model handles frame jitter and slight motion blur effectively during live camera execution.
\end{itemize}

\subsection{Ambient Environmental Lighting and Shadow Variations}
Problems with monocular landmark tracking and AprilTag tracking can occur due to drastic changes in lighting conditions, shadows, and glare \cite{ReviewVirtualKeyboard2020}. Shadows from sunlight or fingers on the paper surface will influence the process of detecting corners.
\begin{itemize}
    \item \textbf{Technical Mitigation:} Adaptively applying a threshold on the images in grayscale before performing quadrilateral extraction in the AprilTag tracker (\texttt{aprilTag/estimater.py}). Moreover, for SVD-based homography estimation, it is required that there should be four clear corners of the tag around the border of the paper surface.
\end{itemize}

\subsection{Anatomical Finger Co-articulation}
There are some intrinsic common tendons in human fingers, specifically those of the ring finger and pinky finger. When the ring finger touches down on a key, the pinky finger is also going to go down slightly, possibly resulting in an accidental touch.
\begin{itemize}
    \item \textbf{Technical Mitigation:} The LSTM model in PyTorch calculates the probability of all five fingers together by applying a multi-label Binary Cross-Entropy loss function ($\mathcal{L}_{\text{BCE}}$). In this way, the neural network is able to distinguish the difference between the desired and accidental touch of keys.
\end{itemize}

\subsection{Absence of Tactile Haptic Surface Feedback}
Touching down on a flat paper surface does not include movement of the key nor the sound from the switch itself. During the first tests, users sometimes felt uncertain about whether their key press was recorded or not, and hence they used additional pressure.
\begin{itemize}
    \item \textbf{Technical Mitigation:} Immediate feedback was delivered to Application 2 Component B when the touch was recognized ($p_k > 0.90$). Specifically, an audible click sound ($< 5\text{ ms}$ latency) and visual feedback by displaying a green square above the touched key on the screen overlay were provided.
\end{itemize}

\section{Self-Reflection}
\label{sec:concl_reflection}
The experience gained from this research is enormous from theoretical and practical points of view concerning computer vision, deep learning, and software engineering.

\subsection{Ideology Regarding Paper-Based Touch Interaction}
Prior to conducting the research, my understanding was that in order to develop a virtual input device, there would be the necessity to use expensive hardware, such as a laser projector, time-of-flight sensor, or electronic glove. Nevertheless, during the research on monocular computer vision and homography geometry, I have realized that there is no necessity of using hardware as the proper software design will be enough to replace the complexity of hardware. For example, using printed paper, ordinary webcam, and a lightweight PyTorch LSTM model on CPU without any GPU, it is possible to develop a virtual input device almost for free. This experience highlighted how software-driven physical interfaces can help make computing input more accessible in low-resource environments.

\subsection{Technical Benefits Gained Across Engineering Domains}
In the process of conducting the project, it was possible to gain practical experience in four fields of computer science:
\begin{enumerate}
    \item \textbf{Practical Computer Vision and Geometry:} Practical experience with calibration of the camera ($K$), rectification of lens distortion ($\mathbf{D}$), AprilTag detection, DLT and SVD algorithms for planar homography ($H \in \mathbb{R}^{3 \times 3}$).
    \item \textbf{Temporal Signals Deep Learning:} Skills of designing and training recurrent PyTorch models, creation of multi-model benchmark pipeline (\texttt{run\_all.py}), using unitless scale normalization ($L_{\text{hand}}$), and preparation of sliding temporal windows ($\mathbf{X}_W \in \mathbb{R}^{5 \times 84}$).
    \item \textbf{Concurrent Multi-Threaded Programming:} Architecture design of producer-consumer threads in PySide6 (\texttt{QThread}, \texttt{QMutex}, \texttt{QWaitCondition}) in order to separate video capturing and neural network inference not to lose frames and maintain throughput ($34.4\text{ FPS}$).
    \item \textbf{Scientific Research and Documentation:} Skills of testing, typing metrics measurement (WPM, CER), chapter organization in scientific research, and IEEE citations in LaTeX documentation.
\end{enumerate}

\subsection{Critical Learning Curves and Methodological Reflections}
\label{subsec:concl_learning_curves}
The most important learning curve that I had to experience with this project is to find a balance between computational capability and the speed of touch events recognition in common consumer laptops. Initially, my approach to the problem was to have a high frame rate ($30\text{--}60\text{ FPS}$) and use complex neural networks, which would require powerful graphics processors (GPUs) and thus make it impossible to use on everyday laptop computers.

However, after profiling and testing of the system, it turned out that the optimal approach is to perform uniform processing at the speed of 12~FPS, which will allow recognizing deceleration of touch events and provide us with an $83.33\text{ ms}$ frame time. This efficient solution allows performing the entire pipeline including vision and deep learning operations without the usage of a graphics processor, and the system will need only $29.09\text{ ms}$ per frame. Unitless hand scale normalization and SVD homography will help to get the same touch events recognition accuracy ($F1 \approx 0.963$) with lower resolution cameras (for example 360p or 480p).

\section{Business Insights and Real-World Application Possibilities}
\label{sec:concl_business}
There are several key benefits of the paper virtual keyboard compared to traditional hardware keyboards. These benefits include an inexpensive manufacturing process, possibility of customization for each user, easy cleaning, and quick deployment across different environments.

\subsection{Sterile Medical Operating Rooms and Clinical Environments}
In medical settings, such as operating rooms and clinics, hardware keyboards and touchscreens are easily contaminated by bacteria because it is difficult to sterilize mechanical keys between patients. However, the paper virtual keyboard with AprilTag anchors can be wiped clean with disinfectant or replaced for just pennies per sheet, providing a clean and sterile input surface for hospital staff.

\subsection{Industrial Cleanrooms and Hazardous Chemical Laboratories}
In semiconductor cleanrooms and chemical laboratories, equipment must be non-sparking and free from dust hazards. Electronic devices can collect particles or cause electrical sparks. The printed paper keyboard with AprilTag arrangement requires no electricity on the table and works completely passively, making it safe in hazardous environments.

\subsection{Low-Cost Educational Classrooms in Developing Regions}
In low-resource computer labs and schools in developing countries, schools often lack financial support to purchase hardware accessories or replace broken keyboards. Using Application 1 Layout Designer (\texttt{designer\_app.py}), teachers can print custom keyboard layouts, math pads, piano keys, or interactive worksheets on standard A4 paper. Connected to a single budget computer and a basic webcam, students can interact with learning software directly on printed paper sheets.

\subsection{Customized Digital Audio Workstations and Accessibility Interfaces}
Hardware macro pads and MIDI controllers for professional software (such as digital audio workstations like Ableton Live, video editing software like Adobe Premiere, or CAD tools) are expensive and have fixed keys. With the help of Application 2, users can switch between different digital action profiles in software (such as recording, editing, or mixing profiles) on the same physical printed sheet without having to print a new sheet. In addition, individuals with motor disabilities can design custom large-button accessibility layouts based on their physical mobility, switching between voice tools, navigation shortcuts, and text typing on the same printed layout.

\section{Future Recommendations}
\label{sec:concl_future}
Although the virtual keyboard performs reliably, there are several promising directions for future research:

\subsection{Integration of Monocular Neural Depth Models}
First, future research could examine integrating lightweight monocular depth estimation models (such as Depth Anything or MiDaS-nano) directly into the pipeline. Combining 2D hand joint coordinates with estimated depth maps could provide millimeter-level Z-axis distance data and improve touch contact classification.

\subsection{Multi-Camera Stereo Vision Fusion}
The use of dual cameras rather than a single camera could help solve the hand occlusion problem. Tracking hand joints from two cameras at different angles would enable 3D triangulation and keep tracking active even when one camera view is partially blocked.

\subsection{Two-Handed (Bimanual) Typing with Dual Window Buffering}
While the current system implementation is designed for single-handed active typing to reduce CPU computation time and avoid occluding fiducial markers on A4 paper sheets, the modular architecture allows scaling to two hands. Future work can configure MediaPipe to track two hands (\texttt{num\_hands=2}) and run separate 5-frame sliding window buffers for each hand. Because the lightweight PyTorch model takes only $2.08\text{ ms}$ per forward pass, dual-hand inference would take approximately $33\text{ ms}$ total on CPU, remaining well within the $83.33\text{ ms}$ 12~FPS frame budget.

\subsection{Spatial Acoustic and Ultrasonic Mid-Air Haptic Feedback}
Ultrasonic mid-air haptic arrays could be used for tactile feedback simulation without mechanical keys. Focusing ultrasonic pressure waves onto the active fingertip during key press confirmation provides physical tactile feedback in the air.

\subsection{Deep Learning Based Fiducial Marker Tracking}
Although the traditional AprilTag 2/3 detector operates efficiently on CPU, future research could evaluate lightweight convolutional marker detectors, such as DeepTag \cite{Zhang2022DeepTag}. Direct estimation of dense keypoints with neural networks can improve homography stability under severe motion blur and finger occlusion.

\subsection{Transformer-Based Temporal Self-Attention Architectures}
Although the uni-directional PyTorch LSTM model performs efficiently ($F1 = 0.963$), future work could evaluate lightweight temporal Transformer models to capture long-range sequential patterns and reduce Character Error Rates during continuous typing.

\section{Chapter Summary}
\label{sec:concl_summary}
This chapter concluded the thesis by reviewing the accomplishment of research objectives, discussing problems and mitigations, reflecting on technical lessons, presenting business use cases, and outlining future research directions.

The research demonstrated that a customizable paper virtual keyboard relying on monocular computer vision, hand scale normalization ($L_{\text{hand}}$), AprilTag planar homography ($H \in \mathbb{R}^{3 \times 3}$), and PyTorch temporal deep learning (\texttt{best\allowbreak\_finger\allowbreak\_touch\allowbreak\_lstm.pth}) achieves reliable, real-time typing performance ($29.09\text{ ms}$ latency, $34.4\text{ FPS}$, $24.6\text{ WPM}$) on standard commodity CPU hardware. By replacing specialized hardware with software engineering, this work provides a flexible, hygienic, and accessible input interface.
